\section{A global complete-intersection formula with actual Green regularity}
\label{sec:global-complete-intersection-green}

The complete-intersection formula of Mihatsch admits the following
explicit globalization when all sections belong to the same line bundle.
The proof is included to check its actual-test scope and its regularity
on the complement of the original cycle.  This is a consequence of known
local formulas, not a claim of a new general solution of strong G.

\begin{proposition}[Two complete intersections of a common line bundle]
\label{prop:global-complete-intersection-strong-green}
Let $X/K$ be smooth projective of pure dimension $n$, and let
$1\le r\le n$ and $q=n-r$.  Let $L$ be a line bundle and let
\[
 s_0=(s_{0,1},\ldots,s_{0,r}),\qquad
 s_1=(s_{1,1},\ldots,s_{1,r})
\]
be two tuples of global sections of $L$.  Suppose that each zero scheme $Z_i$ is purely of codimension $r$
and that the local functions in a trivialization form a regular
sequence at every point of $Z_i$.  Put
$Z=[Z_1]-[Z_0]$, with all generic scheme lengths included, and assume
\[
             |Z|=|Z_0|\cup|Z_1|.                    \tag{CI0}
\]
For example, (CI0) holds if the two cycles have no common irreducible
component.  There is an actual form
\[
       g_Z\in B^{r-1,r-1}(X^{\an}\setminus|Z|)
\]
over the original fixed $K$ such that
\[
       \int_{X^{\an}\setminus|Z|}g_Z\wedge d'd''\eta
                             =-\int_Z\eta
       \qquad\text{for every actual }\eta\in B_c^{q,q}(X^{\an}).
                                                        \tag{CI1}
\]
It defines a continuous functional on full actual target tests of
bidegree $(q+1,q+1)$ for the coefficient-final topology, and an
ordinary CLD current.  The conclusion is an actual off-cycle form,
not merely a current solution.
\end{proposition}

\begin{proof}
Choose a piecewise smooth metric on $L$, and put
$\omega=c_1(L)$ in the actual delta complex.  Such a metric exists
over $K$ without a field extension: choose a sufficiently high power
$H^m$ of an ample bundle so that both $H^m$ and $L\otimes H^m$ are
globally generated, give each the pullback of the standard model
metric on projective space under its section morphism, and take the
quotient metric.  These standard metrics have piecewise smooth
potentials given by finite maxima of logarithms of absolute values;
tensor products and duals preserve the property.  One could instead
use an ordinary smooth metric, whose existence is recalled in
\cite[Section 5.1]{Mih24}.  The argument permits either choice.

For one of the tuples $s=(s_1,\ldots,s_r)$ with zero scheme $Y$, set,
on $U=X^{\an}\setminus|Y|$,
\[
 \varphi=\min_i\{-\log\|s_i\|\},\qquad
 A=\omega-d'd''\varphi,\qquad
 G_s=\varphi\sum_{j=0}^{r-1} A^j\wedge\omega^{r-1-j}.
                                                        \tag{CI2}
\]
The minimum is finite and piecewise smooth on $U$: locally one may
omit sections of strictly smaller norm, and the remaining relevant
ratios are analytic units.  Thus all expressions in (CI2) are actual
forms.  Write $s_i=f_i e$ in a local frame, and put
\[
          \psi=-\log\|e\|,\qquad
          u=\min_i(-\log|f_i|).
\]
Then $\varphi=u+\psi$, $\omega=d'd''\psi$ and $A=-d'd''u$ on $U$.
The morphism $F_s:U\to\mathbf P^{r-1}$ given by the sections identifies
$A$ with the pullback of the standard metric's first Chern form:
the potential in the frame $e$ is $\log\max_i|f_i|=-u$.
In particular $A^r=0$ by dimension and actual pullback.  Both $A$
and $\omega$ are closed under $d'$ and $d''$.  Therefore
\[
 d'd''G_s=(\omega-A)\sum_{j=0}^{r-1}A^j\omega^{r-1-j}
                              =\omega^r\quad\hbox{on }U.       \tag{CI3}
\]

We now check the extension identity against all actual tests; (CI3)
alone would not suffice.  Work on a trivializing analytic open $V$
with compactly supported tests, and keep $q=\dim Y=n-r$.
For $0\le j<r$, the functional denoted by $[A^j]$ is defined on an
actual complementary test $\beta$ by
\[
                 [A^j](\beta)=\int_{V\setminus Y}A^j\wedge\beta.
\]
It is well-defined because both degrees of $\beta$ are $n-j>q$,
so its compact support avoids $Y$ by \cite[Corollary 4.9]{Mih24}.
The first-derivative tests for $[A^j]$ have bidegree
$(n-j-1,n-j)$ or $(n-j,n-j-1)$.  Their maximum degree is still
$n-j>q$, so they also have compact support away from $Y$.
Actual Stokes and the actual closure of $A$ on $U$ consequently give
\[
                         d'[A^j]=d''[A^j]=0
                         \qquad(0\le j<r)              \tag{CI4}
\]
on all these actual tests.

These statements also justify a module operation by a globally defined
actual form: it means wedging that form into the actual test before
evaluating, not multiplying two arbitrary currents.  More explicitly,
for a compact actual test $\beta$ of bidegree
$(n-j-1,n-j-1)$, two applications of Stokes yield
\[
       \int_{V\setminus Y}\psi A^j\wedge d'd''\beta
          =\int_{V\setminus Y} A^j\wedge\omega\wedge\beta.
                                                        \tag{CI5}
\]
The support issue even at $j=r-1$ is worth making explicit.  The first
Stokes integral uses $\psi A^j\wedge d''\beta$; the test
$d''\beta$ has maximum degree $n-j>q$.  The second uses
$A^j\wedge d'\psi\wedge\beta$; the actual form
$d'\psi\wedge\beta$, defined on all of $V$, has the same maximum
degree $n-j>q$.  Thus both expressions have compact support on
$V\setminus Y$.  Likewise $\omega\wedge\beta$ has degrees
$n-j>q$.  Formula (CI5) does not assume that an arbitrary off-cycle
form of bidegree $(r,r)$ has a canonical extension across $Y$.

For $j<r-1$, every compact test of bidegree
$(n-j-1,n-j-1)$ itself avoids $Y$.  Stokes applied directly to
$uA^j$ therefore gives
\[
                    d'd''[uA^j]=-[A^{j+1}]
                                \qquad(0\le j<r-1).     \tag{CI6}
\]
At the top index the degree argument no longer applies.  Instead
\[
                  uA^{r-1}=(-1)^{r-1}u(d'd''u)^{r-1}
\]
is exactly the complete-intersection Green form in
\cite[Theorem 5.5]{Mih24}.  That theorem is explicitly valid for all
actual delta tests, and gives
\[
                       d'd''[uA^{r-1}]=-\delta_Y.        \tag{CI7}
\]
It includes the scheme-theoretic multiplicities of the regular sequence.

Test (CI5)--(CI7), for the $j$th summand, against
$\omega^{r-1-j}\wedge\eta$, where
$\eta\in B_c^{q,q}(V)$.  Since $\omega$ is closed, its exterior
powers commute with the differential in these test identities.
The terms telescope:
\[
 \begin{split}
 d'd''[G_s]
   &= -\delta_Y
      -\sum_{j=0}^{r-2}[A^{j+1}\omega^{r-1-j}]
      +\sum_{j=0}^{r-1}[A^j\omega^{r-j}]\\
   &= -\delta_Y+[\omega^r].
 \end{split}                                                   \tag{CI8}
\]
Every mixed expression in this display is the just-defined functional
obtained by inserting the globally defined power of $\omega$ into an
actual test.  Thus (CI8) is a shorthand for the proved all-actual-test
integral identity and uses no product or pullback of arbitrary currents.
A smooth partition of unity on the compact support of an arbitrary test
glues the local calculation to all of $X$.

Apply (CI8) to the two tuples using the same metric.  Their curvature
terms are the identical actual form $\omega^r$.  Set
$g_Z=G_{s_1}-G_{s_0}$.  Condition (CI0) says exactly that its domain of
actual definition is $X^{\an}\setminus|Z|$.  Subtracting the identities
proves (CI1) with the original cycle, not a cycle having only its class.

Finally, a compact target test of bidegree $(q+1,q+1)$ has support away
from $|Z|$.  On each fixed finite coefficient presentation this
separation is uniform.  These tests need not be top-degree on $X$,
so one uses the degree version of the same argument as Lemma
\ref{lem:uniform-family-support-separation}: near $y\in|Z|$, each
fixed moment map has a sufficiently small affinoid neighborhood whose
tropical image has dimension at most $d(y)\le q<q+1$.
Lemma 4.8(1) of \cite{Mih24} therefore annihilates every test in
that presentation there, independently of its coefficients.
The finite compact localization data and compactness of $|Z|$ give
one common neighborhood of vanishing.  On its compact complement
the fixed actual form $g_Z$ has a finite presentation.  Refining
that presentation with the fixed presentation of the varying test,
wedge product and integration are finite polyhedral coefficient
operations.  Their absolute value is bounded by a finite sum of compact
coefficient seminorms, with fixed finite derivative orders when needed.
This proves continuity on each presentation and hence on the full
coefficient-final space.  The ordinary-current assertion also follows
from \cite[Lemma 5.1]{Mih24}.
\end{proof}

\paragraph{Support and scope.}
The cycles of two regular sections of $L^{\oplus r}$ have the same
Chow class $c_r(L^{\oplus r})\cap[X]=c_1(L)^r\cap[X]$, by the
regular-section formula for the top Chern class (Stacks Project,
Lemma 42.44.1, Tag \texttt{0FA9},
\url{https://stacks.math.columbia.edu/tag/0FA9}).  Their difference is
therefore rationally, and hence geometrically homologically, trivial.
That class equality is only a verification of the hypothesis class:
the formula (CI1) is for the original difference of schemes with its
multiplicities.  If a common component cancels completely so that
(CI0) fails, the same calculation yields a form only on the complement
of the union.  Extension across the canceled component has not been
proved and is not asserted as a solution of the original-support problem.
Finite rational linear combinations of the proved differences give the
same conclusion only when their combined prescribed support is still
the union of the supports on which the summand forms are defined.

In particular, taking $L=\pi_2^*\mathcal O_{\mathbf P^r}(1)$ on
$\mathbf P^q\times\mathbf P^r$ recovers actual Green forms for the two
fibers used in the nonrepresentation example.  The specific family
functional there has no actual representative off the cycle when
$r\ge2$, while this different solution does.  Their difference is a
$d'd''$-closed current; this illustrates the regularization issue without
solving it for arbitrary families or arbitrary homologically trivial
cycles.

\paragraph{The resulting limited Arakelov pairing.}
For the specified actual $g_Z$ and an original rational algebraic cycle
$W$ of dimension $r-1$ disjoint from $|Z|$, the number
\[
                       h_{g_Z}(Z,W)=\int_{W^{\an}}g_Z|_W
\]
is well-defined by actual pullback and top-degree integration, with
the original multiplicities of $W$. Properness of $W$ gives compact
support. The actual projection and finite base-change formulas proved
earlier apply, after choosing a cutoff equal to one near $W$ if needed.
This is a local height contribution for the chosen Green object; no
independence of the choice of metric, section tuple, or Green form is
asserted. It does not define products of arbitrary Green currents or
a general height pairing for cycles meeting $|Z|$.
